\section{Brianchon's theorem}

\subsection{Statement and strategy}

Let a hexagon $V_1\ldots V_6$ be circumscribed about a nondegenerate projective
conic: its six sides are tangent to the conic.  Then the three diagonals
joining opposite vertices,
\[
 V_1V_4,\qquad V_2V_5,\qquad V_3V_6,
\]
are concurrent.  We prove this as the exact point--line dual of Pascal's
theorem, but we also give the polynomial certificate for Pascal so that no
unproved invocation of a general duality principle remains.

\subsection{Pascal on a standard conic}

Work first with the parametrized conic
\[
 p(t)=(t^2,t,1),
 \qquad xz-y^2=0.
\]
The secant through $p(s)$ and $p(t)$ has line covector
\begin{equation}\label{eq:secant-line}
 L(s,t)=(1,-(s+t),st),
\end{equation}
up to the harmless scalar $s-t$.  For six parameters $a,b,c,d,e,f$, define
the three Pascal points
\[
 X=L(a,b)\meet L(d,e),
 \quad Y=L(b,c)\meet L(e,f),
 \quad Z=L(c,d)\meet L(f,a).
\]
Direct cross products give
\begin{align*}
 X={}&\bigl(ab(d+e)-de(a+b),\ ab-de,\ a+b-d-e\bigr),\\
 Y={}&\bigl(bc(e+f)-ef(b+c),\ bc-ef,\ b+c-e-f\bigr),\\
 Z={}&\bigl(cd(f+a)-fa(c+d),\ cd-fa,\ c+d-f-a\bigr).
\end{align*}
The complete Pascal certificate is the vector identity
\begin{equation}\label{eq:pascal-certificate}
 (c-f)X+(d-a)Y+(e-b)Z=0.
\end{equation}
Every coordinate of the left side expands to zero.  Hence $X,Y,Z$ are
linearly dependent and therefore collinear.  Repeated parameters are allowed
whenever the limiting secant is interpreted as the tangent; alternatively,
for the present Brianchon application we use six tangent lines directly and
need no limiting argument.

\subsection{The dual conic of tangent lines}

The tangent to $xz-y^2=0$ at $p(t)$ is
\begin{equation}\label{eq:tangent-line}
 \tau(t)=(1,-2t,t^2),
\end{equation}
because the gradient at $(t^2,t,1)$ is $(1,-2t,t^2)$.  Regard a line covector
as a point of the dual projective plane.  The six tangents
\[
 \ell_1=\tau(a),\ \ell_2=\tau(b),\ldots,\ell_6=\tau(f)
\]
are six points on a dual conic, projectively equivalent to the conic used in
the Pascal calculation.

The vertices of the circumscribed hexagon are
\[
 V_i=\ell_i\meet\ell_{i+1}
 \quad(i\bmod 6).
\]
The first opposite-vertex diagonal has covector
\begin{align*}
 d_1
  &=V_1\join V_4\\
  &=(\ell_1\cross\ell_2)\cross(\ell_4\cross\ell_5).
\end{align*}
In the dual plane this same triple is the intersection of the dual secants
$\ell_1\ell_2$ and $\ell_4\ell_5$, i.e. the first Pascal point.  Similarly,
\[
 d_2=V_2\join V_5,
 \qquad d_3=V_3\join V_6
\]
are the second and third Pascal points.  Therefore
\[
 \det(d_1,d_2,d_3)=0
\]
by \eqref{eq:pascal-certificate}.  But a determinant of three line covectors
vanishes exactly when the three lines are concurrent.  If $d_1,d_2$ are
distinct, the explicit concurrency point is
\[
 P=d_1\meet d_2=d_1\cross d_2,
\]
and $d_3\cdot P=\det(d_3,d_1,d_2)=0$.

\subsection{Reduction of a general conic}

Represent a nondegenerate rational conic by an invertible symmetric matrix
$S$, after multiplying the equation by $2$ to remove half-coefficients.  A
rational tangent covector $\ell$ is a point of the dual conic, and its contact
point is represented by
\[
 p=\operatorname{adj}(S)\ell.
\]
Indeed $Sp=\det(S)\ell$, so the tangent at $p$ is $\ell$; the dual-conic
equation ensures that $p$ lies on the original conic.  Thus one of the six
supplied tangents gives a rational contact point without solving a quadratic.

Using bounded search through the coordinates and $2\times2$ minors, choose a
projective frame that sends this contact point to $(0:0:1)$ and its tangent to
$x=0$.  In that frame the conic equation has the form
\[
 xz+A x^2+Bxy+C y^2=0.
\]
The coefficient of $xz$ is nonzero because $x=0$ is the tangent, and $C\ne0$
because the conic is nondegenerate.  The explicit substitutions
\[
 z'=z+Ax+By,
 \qquad z''=-z'/C
\]
carry the equation to $xz''-y^2=0$.  All transformations are rational
adjugate-and-rescaling operations under named nonzero minors; no square root
or choice principle is involved.

The tangent covectors of the standard conic have homogeneous parameterization
\[
 \tau(s:t)=(s^2,-2st,t^2).
\]
If a finite-parameter chart is desired for the polynomial identity
\eqref{eq:pascal-certificate}, choose a rational M\"obius change of parameter
whose point at infinity is not one of the six supplied parameters.  Seven
fixed candidates suffice: each of the six points forbids at most one
candidate.  This is a bounded decidable search.  Alternatively, homogenize
the three coordinates of \eqref{eq:pascal-certificate} in the six parameter
pairs.

Projective transformations preserve incidence and multiply cross products and
determinants only by explicit nonzero scalar factors.  Applying the standard
certificate in the transformed dual plane and transporting it back proves
Brianchon's theorem for the original conic.  For Lean, a useful intermediate
interface takes the normalizing matrix and six homogeneous parameters as a
certificate; a separate lemma implements the intrinsic construction above.

\subsection{Constructive content}

Brianchon's theorem uses no topological limiting argument and no appeal to a
semantic ``principle of duality.''  The dualization is the syntactic exchange
of point triples and line triples, while join and meet are both the same cross
product.  The proof is a fixed polynomial identity followed by the explicit
witness $d_1\cross d_2$ under a nonparallelity hypothesis.
